\section{Evaluation Protocol}
\label{appendix:protocol}


\subsection{Initial Conditions and Success Criteria}
\label{sec:ic_design}

\textbf{Hardware trials.}
We use a stratified initial condition (IC) design to ensure
fair comparison across controllers with different constraint-handling
capabilities. Strata are defined mainly by the cart's starting
position (available maneuvering space), and at the rail by the targeted
ball side. The cart is positioned automatically via closed-loop
velocity control (P-controller, 10~mm tolerance), and the ball is
displaced from center via brief cart jolts.

The 50-trial schedule comprises five strata:
\begin{itemize}
    \item \textbf{Center (20\%, 10 trials):} Cart at $\pm 0.10$~m,
          ball jolted off center.
    \item \textbf{Mid-rail (32\%, 16 trials):} Cart at $\pm 0.40$~m,
          ball jolted off center. Tests recovery
          with asymmetric rail budget ($\sim$0.38~m one side,
          $\sim$1.18~m other).
    \item \textbf{Near-wall (28\%, 14 trials):} Cart at $\pm 0.65$~m,
          ball jolted off center. Tests recovery
          under tight constraints (only 0.13~m clearance on one side).
    \item \textbf{At-wall-opposite (10\%, 5 trials):} Cart at
          $\pm 0.78$~m (workspace limit, axis units), ball displaced to
          \emph{opposite} side. Feasible for all controllers;
          tests recovery speed.
    \item \textbf{At-wall-same (10\%, 5 trials):} Cart at
          $\pm 0.78$~m (workspace limit, axis units), ball displaced to
          \emph{same} side as cart. Reveals structural limitations
          of controllers without explicit constraint handling.
\end{itemize}

Trials alternate left/right within each stratum and are shuffled
deterministically (seed~42) so every controller sees the identical
sequence.  The schedule is logged alongside results for reproducibility.

Beyond the ball-side selection in the two at-wall strata, the ball's
initial angle $\theta_0$ is not controlled: it arises from
the brief cart jolt used to displace the ball. Empirically,
$\theta_0 \in [-0.091, 0.081]$~rad across the benchmark trials
(mean $|\theta_0| \approx 0.062$~rad, standard deviation $\approx 0.019$~rad),
and it stays in the same range in every stratum, with per-stratum mean
$|\theta_0|$ ranging only from $0.059$~rad (at-wall-opposite) to $0.071$~rad
(at-wall-same). The initial ball displacement is therefore comparable across
strata, so stratum difficulty reflects maneuvering room rather than initial
ball displacement.

The jolt also places the ball outside the central dip in nearly every trial.
With the surface parameters of Section~\ref{appendix:dynamics}, the dip's
restoring region reaches only to $|\theta| \approx 0.0115$~rad. Released at rest
inside that angle with the cart held still, the ball returns to center, and
released outside it, the ball rolls away down the convex arc. All but five of the 650 benchmark trials start beyond that angle, at a
mean $|\theta_0|$ some five times larger, so a controller has to carry the ball
back across the destabilizing part of the arc before the dip can capture it.

One caveat applies to the achieved cart start in the at-wall strata, where the
cart is driven toward the rail ($|x_0| = 0.777$~m) before the ball is displaced.
In the at-wall-\emph{same} stratum the cart reached the physical limit on about
$83\%$ of the 13 benchmark controllers' trials, but this was not uniform across
controllers: most reached the rail on all five trials, whereas for a few,
notably NMPC and TD3, on three of their five trials the automated placement
left the cart near center ($|x_0| \approx 0.08$~m) rather than at the rail, a
materially easier start. Both still succeeded on all five, and within this stratum the
success rate does not differ significantly between trials that reached the limit
and those that did not ($50/54$ versus $9/11$; Fisher's exact test $p = 0.27$),
so the effect on the reported comparison is small. The at-wall-\emph{opposite}
stratum is less uniform still: the cart reached the limit for only some
controllers and started nearer center ($|x_0| \approx 0.09$~m) for the others.
We therefore read both at-wall strata as recovery-speed checks for the
controllers that did not reach the rail, rather than strictly equal-condition
cross-controller robustness comparisons.

\textbf{Simulation controller benchmarks} (the RL and classical
comparisons run in simulation) use the same 50-trial stratified schedule
as the hardware protocol (cart at
$\pm 0.10, \pm 0.40, \pm 0.65, \pm 0.78$~m). The simulator also provides a
legacy uniform-sampling mode,
\begin{align}
  x_0 &\sim \mathcal{U}(-0.75,\;0.75)~\text{m}, \nonumber\\
  \theta_0 &\sim \mathcal{U}(-0.074,\;0.074)~\text{rad}, \nonumber\\
  \dot{x}_0 &= \dot{\theta}_0 = 0,
\end{align}
but it was not used for the reported benchmarks.

\textbf{Success criterion.}
A trial is \textbf{successful} if the ball angle satisfies
\begin{equation}
  |\theta(t)| < 0.01~\text{rad}
\end{equation}
\emph{continuously} for at least $1.0$~s.
The \textbf{settling time} $t_s$ is defined as the wall-clock time at
which the ball \emph{enters} this band and remains within it for the
required duration (i.e.\ the entry instant, not the end of the window).
The maximum trial duration is $30$~s; trials not meeting the settling
criterion within this window are recorded as failures.

\subsection{Performance Metrics}

We evaluate controllers using the following metrics, computed with
SciPy~\cite{2020SciPy-NMeth} and NumPy~\cite{harris2020array}:

\begin{itemize}
    \item \textbf{Settling time} ($t_s$): wall-clock time from trial start
    until the success criterion is first met (as defined above). Medians,
    interquartile ranges, and the settling significance tests are computed
    over \emph{successful trials only}, matching the main paper's Table~II
    statistics; success rate is reported separately, so reliability is not
    folded into the settling comparison.

    \item \textbf{Success rate}: fraction of trials with $t_s < 30.0$~s,
    with 95\% Wilson score intervals. Wilson intervals replace a bootstrap
    because a perfect $50/50$ record bootstraps to a degenerate
    $[100, 100]$ interval, understating uncertainty at the boundary.

    \item \textbf{Control effort}: root-mean-square commanded cart
    velocity in m/s, saturated at the $\pm 0.9$~m/s axis limit the
    actuator enforces, computed per trial and averaged over all trials.

    \item \textbf{Significance testing.}
    Pairwise Mann-Whitney $U$ rank tests
    (\texttt{scipy.stats.mannwhitneyu}) on the successful-trial
    settling distributions, Bonferroni-corrected~\cite{bonferroni1936teoria}
    across all $\binom{13}{2}=78$ pairs ($\alpha_c = 0.05/78$). A rank test
    is used because the settling distributions are right-skewed, so a
    mean-based $t$-test is inappropriate. Success-rate differences use
    Fisher's exact test (\texttt{scipy.stats.fisher\_exact}) over the same
    family. Control-effort differences are tested the same way, as a separate
Mann-Whitney family over the per-trial RMS effort across all $78$ pairs, with
the full pairwise matrix written to \texttt{exp1\_effort\_significance.csv}.
Effect sizes are reported as Cohen's $d$
    ($d = (\bar{x}_i - \bar{x}_j)/s_{\mathrm{pooled}}$,
    $s_{\mathrm{pooled}} = \sqrt{(s_i^2 + s_j^2)/2}$, equal-weight pooling)
    to complement the $p$-values. Because the 78 pairwise tests are computed
    on overlapping controller data rather than independent samples, the
    Bonferroni correction is conservative here, so the number of pairs flagged
    significant is best read as a lower bound. All statistics are reproduced by
    \texttt{paper/scripts/compute\_statistics.py}.
\end{itemize}

\subsection{Experiment 1: Baseline Performance}

\paragraph{Objective:}
Establish the nominal performance hierarchy under realistic but
unperturbed sensing conditions, using difficulty-stratified ICs that
test regulation, recovery, and constraint handling.

\paragraph{Protocol:}
We evaluate each controller over 50 trials on the physical
platform at the nominal control frequency of 20~Hz
($T_s = 0.05$~s).
Sensor readings use the 34.5~ms timing-budget sequential ToF
configuration (Appendix~\ref{appendix:sensor_characterization}) with inherent noise
($\sigma_{\text{ToF}} \approx 1.6$--$1.7$~mm at the arc center);
no additional artificial noise is injected.

\paragraph{Rationale:}
The stratified IC design ensures that:
(a) no controller is unfairly penalized by impossible starting
conditions;
(b) center strata establish baseline regulation performance;
(c) near-wall/at-wall strata discriminate controllers by their
constraint-handling and recovery capabilities;
(d) the identical IC sequence across all controllers removes
IC-related confounds from pairwise comparisons, with the exception of
the achieved cart start in the two at-wall strata, which was not uniform
across controllers (Section~\ref{sec:ic_design}).
